\section{Additive Manufacturing Routes for HEAs}
\label{sec:am}

AM of HEAs has been reviewed comprehensively by Cagirici et~al. \cite{cagirici2024} and Han et~al. \cite{han2025}; here we summarize the process--property consequences that matter for marine components, then discuss feedstock, which is the binding constraint for marine-scale deployment. Table~\ref{tab:am_processes} (Section~\ref{sec:tabfig}) compares the routes.

\subsection{Powder bed fusion: LPBF and SEBM}

Laser powder bed fusion (LPBF) melts thin powder layers with a focused laser, producing the highest cooling rates ($\sim10^{5}$--$10^{6}$~K/s) and the finest as-built microstructures of any practical HEA route: cellular and dendritic substructures with sub-micron dislocation cells and strong crystallographic texture along the build direction \cite{cagirici2024,han2025}. The same rapid solidification that refines the structure is what suppresses elemental segregation in multiphase HEAs; this is the microstructural origin of the improved chloride corrosion resistance reported for LPBF CoCrFeMnNi and SLM AlCoCrFeNi$_{2.1}$ relative to cast material \cite{lpbf_cantor_2024,slm_ehea_2025}. The penalties are limited build volume, slow build speed, and susceptibility to lack-of-fusion and keyhole porosity when parameters are off-window; high-Al compositions additionally face solidification cracking \cite{cagirici2024}.

Selective electron beam melting (SEBM/EBAM-class processes) operates under vacuum with a heated powder bed, reducing thermal gradients and residual stress. This makes SEBM the preferred powder-bed route for crack-prone and refractory systems \cite{abdullah2024rhea}, and its preheat capability is valuable for BCC HEAs with ductile-to-brittle transitions. Surface finish and resolution are coarser than LPBF, so wetted surfaces intended for corrosion service require machining or polishing.

\subsection{Directed energy deposition and laser metal deposition}

DED/LMD feeds powder or wire into a moving melt pool. Its marine significance is twofold. First, it enables \emph{cladding and repair}: HEA overlays on marine steel substrates or worn bronze/steel components deliver corrosion- and wear-resistant surfaces at a fraction of a monolithic HEA part's cost; laser-clad AlCoCr$_x$FeNi coatings on steel have been studied specifically in this light \cite{jiang2019clad}. Second, DED supports \emph{in-situ alloying and compositionally graded deposits}, which researchers have used both to screen compositions and to functionally grade from a steel substrate to a corrosion-resistant HEA surface \cite{cagirici2024}. Within this route, Ti- or Cu-doped AlCoCrFeNi deposits and Nb--Mo--Ta--W compositions designed for marine service have been reported; the former showed improved corrosion resistance, while the latter---despite better overall performance than 316L---exhibited localized pitting associated with minor compositional segregation, illustrating the defect sensitivity discussed in Section~\ref{sec:corrosion} \cite{cagirici2024}. Control of dilution and interfacial galvanic behavior is the key engineering issue for clad components.

\subsection{Wire arc additive manufacturing}

WAAM deposits material by arc welding of wire feedstock at deposition rates and build volumes far beyond powder-bed processes, which is why it is the route of record for large marine and offshore structures \cite{raultaiwade2021}. For HEAs the bottleneck has been wire availability, and three strategies have emerged: cable-type twisted wires (demonstrated for MoNbTaWTi, yielding single-phase BCC deposits with hardness near 533~HV and compressive strength stable between 500 and 1000~$^{\circ}$C \cite{liu2021waam}), combined-cable wire for Al--Co--Cr--Fe--Ni \cite{shen2021ccw}, and metal powder-cored wires, now demonstrated for Cantor-type CoCrFeMnNi walls with homogeneous element distribution and competitive mechanical properties \cite{andersson2026waam}. WAAM HEA deposits are coarse and columnar, with cyclic reheating from successive passes; Mn evaporation, oxide and sulfide inclusions, and mild elemental segregation (notably Mn) are recurring metallurgical issues \cite{andersson2026waam,raultaiwade2021}. Because WAAM parts are typically machined to finish, the as-deposited corrosion performance is less critical than the properties of the final machined surface---but inclusions and lack-of-fusion defects remain potential pit initiation sites.

\subsection{Binder jetting and other routes}

Binder jetting followed by sintering offers high build speed and low cost but leaves residual porosity unless infiltration or hot isostatic pressing (HIP) is applied; porous binder-jet HEA structures have been proposed for lightweight or functional applications \cite{cagirici2024}. For load-bearing, seawater-wetted marine parts, residual porosity is a corrosion liability (Section~\ref{sec:corrosion}), so binder jetting is currently a secondary route for this application domain.

\subsection{Feedstock: the binding constraint}

Marine deployment is ultimately limited less by process physics than by feedstock economics. Gas-atomized pre-alloyed HEA powders are expensive, and the critical-raw-material content (Co, Ni, Ta, Nb) of many attractive compositions raises both cost and supply-risk questions \cite{cagirici2024,han2025}. Wire production for WAAM---drawing, cabling, or powder-coring---is immature for most HEA compositions, and element evaporation (Mn, Cr) during wire manufacture and deposition must be compensated in the chemistry \cite{andersson2026waam}. Powder reuse and oxidation control in the humid, saline atmosphere of shipyard environments add further qualification requirements. A credible marine HEA program therefore starts with a feedstock strategy, not a printer.
